\section{File System Consistency and Recovery}
\label{sec:recovery}

\subsection{What a Crash Does}

Creating a file is not one disk write but a sequence of them. When everything is normal, the system call carries out the following steps (\Cref{fig:createsteps}):
\begin{enumerate}
  \item read the free-space management data from the disk;
  \item update the free-space management data on the disk;
  \item read the file allocation information from the disk;
  \item update the file allocation information and write it to the disk;
  \item read the directory entry from the disk;
  \item update the directory entry on the disk;
  \item write the data to the allocated blocks.
\end{enumerate}

\begin{figure}[htbp]
\centering
\begin{tikzpicture}[stp/.style={draw, fill=rowc, font=\scriptsize, minimum width=58mm, minimum height=5.5mm, inner sep=2pt},
                    note/.style={font=\scriptsize, align=left, anchor=west}]
  \foreach \t [count=\i] in {Read free-space management data from disk.,
                            Update free-space management data on disk.,
                            Read file allocation information from disk.,
                            Update file allocation information on disk.,
                            Read directory entry from disk.,
                            Update directory entry on disk.,
                            Write data to allocated blocks.} {
    \node[stp] (s\i) at (0,-0.75*\i) {\i.\ \t};
  }
  \foreach \i [evaluate=\i as \j using int(\i+1)] in {1,...,6} \draw[-{Stealth[length=1.3mm]}] (s\i) -- (s\j);
  \draw[sigc, thick, decorate, decoration={brace, amplitude=4pt}] (3.1,-1.2) -- (3.1,-1.8)
    node[note, midway, xshift=4pt, text=black] {Crash after 1: no problem.};
  \draw[sigc, thick, decorate, decoration={brace, amplitude=4pt}] (3.1,-1.95) -- (3.1,-3.3)
    node[note, midway, xshift=4pt, text=black] {Crash after 2 or 3: free space is allocated to nothing.\\ \emph{Error 1}: the FS loses free space.};
  \draw[sigc, thick, decorate, decoration={brace, amplitude=4pt}] (3.1,-3.45) -- (3.1,-4.8)
    node[note, midway, xshift=4pt, text=black] {Crash after 4 or 5: allocated blocks are chained together\\ but belong to no file. \emph{Error 2}: free space lost too.};
  \draw[sigc, thick, decorate, decoration={brace, amplitude=4pt}] (3.1,-4.95) -- (3.1,-5.6)
    node[note, midway, xshift=4pt, text=black] {Crash after 6: the entry points to uninitialized data.\\ \emph{Error 3}: the content of the file is wrong.};
\end{tikzpicture}
\caption{The disk updates of a file creation, and the inconsistency left by a crash after each one (after slides 216--222).}
\label{fig:createsteps}
\end{figure}

After the OS comes back, can you \emph{detect} these inconsistencies? Can you \emph{recover} from them? Sometimes you can detect an error but have no way to recover from it. The tools for file system checking and recovery are \code{fsck} on Linux and macOS, and \code{chkdsk} and \code{scandisk} on Windows.

\subsection{Error 1: Free Space Allocated to Nothing}

To detect it, check the free-space management data against the entire set of allocated blocks. \code{fsck} can do this for ext2/3/4:
\begin{enumerate}
  \item construct a bitmap of the allocated blocks from all the inodes in the inode table, by following every block pointer;
  \item compare it with the free-space bitmap.
\end{enumerate}
A block that is marked in use in the free-space bitmap but belongs to no inode is a \emph{missing block}. To recover, reset its bit in the free-space bitmap, and reset fields in the superblock if necessary (the count of free blocks).

\subsection{Error 2: Allocated Blocks That Belong to No File}

Here an inode has been written and points to blocks, but no directory entry refers to the inode, so checking the inode table is not enough: the whole directory tree must be checked.
\begin{enumerate}
  \item Traverse the entire directory tree, starting from the root.
  \item Locate the allocated inodes from the inode number in every directory entry, and record them in an inode allocation table.
  \item Construct the block allocation bitmap from the inodes discovered this way.
  \item (a) Compare this block bitmap with the free-space bitmap; (b) compare the inode allocation table built from the tree with the inode bitmap.
\end{enumerate}
This discovers missing blocks, as before, and \emph{dangling inodes}: inodes that are allocated but referred to by no directory entry. To recover, reset the free-space bitmap and the superblock fields as before, and delete the dangling inodes, much as a garbage collector frees unreachable objects. (In practice, \code{e2fsck} offers to reconnect such inodes into the directory \code{/lost+found} instead, so that their data can be inspected.)

\subsection{Error 3: Journaling}

The third error cannot be detected this way: every structure is consistent with every other one, only the data is wrong. Its solution is a \emph{journaling file system}.

\begin{example}[The Log Book]
  Your boss orders you to do a set of tasks, and you write them all down in a log book:
  \begin{enumerate}
    \item Go to GameStop and buy the boss GTA6.
    \item Pick up the boss's girlfriend.
    \item Drive her back to her home.
    \item Buy the boss a green cap on the way back.
  \end{enumerate}
  You cross out each task when it is completed. Unfortunately, a car accident happens and you lose all your memory. Your boss sends a colleague to finish your job; the colleague knows nothing about your progress, and the boss has forgotten what he asked of you. The log book comes in handy: it lists exactly the tasks not yet crossed out.

  In file system terms, the boss is the user program, the tasks are the FS operations it invokes, the accident is a system crash in which all memory is lost, and the colleague is the file system recovery tool. The log book is the journal.
\end{example}

\begin{definition}[Journal and Transaction]
  A \emph{journal} is the log book of a file system, kept inside the file system, i.e., on the disk. A set of file system operations is grouped into an \emph{atomic transaction}: either all of its operations complete successfully, or none does (recall test-and-set in \Cref{sec:cs}). A transaction records all the changes that will be made to the FS, and every transaction is written to the journal.
\end{definition}

\begin{keypoint}[Write-Ahead Principle]
  All changes to the file system are written to the journal first: the changes to the metadata (inodes, directory entries, \dots) and the file content. Then the FS-related system call returns to the user process. Meanwhile, the entries in the journal are \emph{replayed}: the kernel carries out the recorded actions on the actual file system, marks each change as done when it is committed, and removes the transaction from the journal when all of its changes are committed.
\end{keypoint}

\begin{figure}[htbp]
\centering
\begin{tikzpicture}[stp/.style={draw, fill=rowc, font=\scriptsize, minimum width=52mm, minimum height=5.5mm, inner sep=2pt}]
  \foreach \t [count=\i] in {Update free-space management data on disk.,
                            Update file allocation information on disk.,
                            Update directory entry on disk.,
                            Write data to allocated blocks.} {
    \node[stp] (s\i) at (0,-0.8*\i) {\t};
  }
  \foreach \i [evaluate=\i as \j using int(\i+1)] in {1,2,3} \draw[-{Stealth[length=1.3mm]}] (s\i) -- (s\j);
  \begin{scope}[on background layer]
    \node[fit=(s1)(s4), fill=freec, inner sep=5pt] (tr) {};
    \node[fit=(tr), fill=metac, inner sep=6pt, label={[font=\scriptsize]below:a transaction inside the journal}] {};
  \end{scope}
  \draw[sigc, thick, {Stealth[length=2mm]}-] (s1.east) -- ++(0.9,0) node[right, font=\scriptsize, draw=sigc, fill=white, inner sep=2pt] {current task pointer};
  \node[draw, fill=greybg, cylinder, shape border rotate=90, aspect=0.25, minimum width=13mm, minimum height=11mm, font=\scriptsize] (disk) at (6.3,-3.0) {Disk};
  \draw[arr] (s1.east) ++(0,-0.1) to[out=-10, in=150] node[above right, font=\scriptsize, pos=0.55] {commit changes} (disk.north west);
  \node[font=\scriptsize, align=left, anchor=west] at (3.4,-4.2) {Metadata updates (1--3) and the\\ file content update (4).};
\end{tikzpicture}
\caption{A file creation as one transaction in the journal. The kernel commits its actions to the actual file system one by one, then removes the transaction (after slides 237--239).}
\label{fig:journal}
\end{figure}

\subsubsection{Recovery with a journal}
The journal is extremely useful when the system crashes. After a crash, the file system should recover itself, and with the journal, recovery means \emph{replaying the uncommitted transactions}. The recovery tool finds each outstanding transaction in the journal and replays its uncommitted operations; updates that were finished before the crash need no replay.

What if the crash happens while the system is writing the journal itself? Then, during recovery, the transaction is found to be \emph{incomplete}: its last operation was not fully recorded, so the transaction cannot be completed. Since a transaction is atomic, the operations of the incomplete transaction that were already carried out (if any) are \emph{rolled back}, and the transaction is removed. Either way, after recovery the FS is in a state where each transaction happened completely or not at all.

\subsection{Speed versus Reliability}

With full journaling, every piece of data is written twice: once in the journal and once in the actual file system. ext4 therefore offers three journaling modes.
\begin{itemize}
  \item \emph{Journal mode} (slowest, most reliable): write both the metadata and the file content to the journal first, then write them to the file system.
  \item \emph{Ordered mode} (the default; pretty good speed and reliability): flush (force the writing of) the file content to the file system first, then write the metadata to the journal. File content is not journaled, but metadata never points to data that was not written, so Error 3 cannot occur.
  \item \emph{Write-back mode} (fastest, but incorrect data may appear): the file content may be flushed to the file system \emph{after} the metadata has been written to the journal, so after a crash a file may contain stale data.
\end{itemize}
